\input{preamble}
\usepackage{tikz}
\setlength{\emergencystretch}{1em}
\input{paper2/results/evidence_status}
\usetikzlibrary{arrows.meta,positioning,shapes.geometric,fit,calc}
\definecolor{scenarioBlue}{HTML}{3269A8}
\definecolor{encoderPurple}{HTML}{7954A1}
\definecolor{graphGreen}{HTML}{348F70}
\definecolor{ebmOrange}{HTML}{D68B23}
\definecolor{samplerTeal}{HTML}{298F99}
\definecolor{milpRed}{HTML}{B54D52}
\date{12 November 2026}
\title{Which Optimization Method Delivers the Required Energy Service?\\
Deadline-Constrained Comparisons and Limits of Learned Proposals}
\begin{document}
\maketitle
\begin{abstract}
Energy-system optimization must return an actionable dispatch before a deadline,
but physical validity, economic quality and continuity of service are distinct
requirements. We compare direct HiGHS optimization, rolling-horizon proposals
with repair, a single learned latent proposal, and an energy-guided portfolio
of sixteen proposals on 600 constructed scenarios. Independent execution
budgets include all online computation and process closure. At 120 seconds,
HiGHS delivers 600 valid outputs and 505 within five percent of a common
reference, whereas rolling horizon preserves service without unserved energy
on 577 scenarios, compared with 482 for HiGHS. Requiring both economic quality
and continuity yields 475 successes for HiGHS, 248 for the portfolio, 222 for
the single learned proposal and 76 for rolling horizon. Paired comparisons
identify narrower opportunities for learned proposals, including useful
short-budget coverage on renewable/storage scenarios. They do not establish
an incremental benefit of the learned energy: candidate count and repair
remain competing explanations. We specify independent, matched-ablation and
structural-transfer evaluations, including a separately trained variable-horizon
model and an external solver. \pIIExtensionStatus
The supported conclusion is conditional: the useful method depends
on the requested service and allocated time, while generalization and
prospective selection require additional evidence.
\end{abstract}

\section{Introduction}
Operational energy decisions coordinate generation, storage, flexible demand
and exchanges under intertemporal constraints. A dispatch delivered after its
decision deadline is unavailable when needed. A dispatch delivered on time
may still shed load, and one preserving all demand may have a substantial
economic penalty. Comparing optimization methods therefore requires a
specified service requirement, rather than a single unconstrained ranking.

Mixed-integer linear programming (MILP) represents commitment and operating
modes jointly with continuous energy flows. Learning can accelerate different
parts of this process \citep{bengio2021}. Graph-based branching learns decisions
inside an exact search \citep{gasse2019}; Neural Diving generates partial
assignments completed by a solver \citep{nair2021}; learned large-neighborhood
search chooses subsequent optimization neighborhoods \citep{sonnerat2022}.
In unit commitment, historical solutions can support constraint screening,
initial solutions and restricted search \citep{xavier2021}. Recent work combines
graph-based commitment prediction, feasibility restoration and large-neighborhood
search \citep{yang2025lns}. These approaches establish the relevance of
learning-assisted optimization, but their guarantees and computational budgets
are not interchangeable.

Our learned methods propose zonal operating strategies in a latent space,
then reconstruct a dispatch through a common repair cascade. Their operational
value depends on the complete chain: scenario encoding, proposal quality,
repair effort, validation and termination. Comparing one proposal with sixteen
energy-guided proposals confounds guidance with search breadth. Likewise,
feasibility after substantial repair does not demonstrate that the proposal
itself generalized. Fair benchmarking must expose these distinctions and
retain unsuccessful executions \citep{beiranvand2017}.

A graph representation can accept variable-size inputs while its learned
mapping fails on changed local structures \citep{yehudai2021}. The complete
optimization stack may also contain fixed index embeddings or horizon limits.
New draws, structural transformations and horizons absent from training
consequently answer different questions; none follows from architectural
compatibility alone.

We ask: \emph{under the same operational deadline, which methods deliver a
physically valid, economically acceptable dispatch that preserves service,
and how robust are their useful regimes to changes in the scenario?}
The study has three components: a failure-aware map of time--cost--service
trade-offs; a matched design isolating energy guidance; and separate tests of
frozen-model transfer and a newly trained variable-horizon chain. The first
component is supported by a completed campaign. \pIIIntroStatus
This separation prevents
planned controls from being mistaken for demonstrated contributions.

\section{Methods}
\subsection{Scenario and physical optimization problem}
One scenario specifies zones $\mathcal Z$, periods $\mathcal T$, hourly duration
$\Delta t$, a capacitated transport network, demand, renewable availability,
installed capacities, operating limits, initial and terminal states, and costs.
These inputs remain fixed across methods and deadlines. Different deadlines
create different executions of the same physical instance.

The MILP coordinates aggregate zonal thermal commitment, nuclear generation,
renewable curtailment, batteries, reservoir and pumped hydro, demand response,
line flows and external trade. At every zone and period, generation, storage
flows, exchanges and flexible demand satisfy nodal balance, with explicitly
penalized unserved demand and excess generation. The objective is
\begin{equation}
 C=\sum_{z,t}\Delta t\bigl(c^{\rm gen}_{zt}+c^{\rm flex}_{zt}
       +c^{\rm cyc}_{zt}+c^{\rm spill}_{zt}+v_zP^{\rm uns}_{zt}\bigr)
       +\sum_{z,t}c_z^{\rm start}v_{zt}
       +C^{\rm import}-R^{\rm export}.
 \label{eq:cost}
\end{equation}
The parenthesized terms are cost rates; start-up costs are charged per event.
Export revenues can make net costs negative. Thermal transitions obey
$u_t-u_{t-1}=v_t-w_t$, with mutually exclusive start-up and shutdown events and
explicit ramp limits. Storage tracks efficiency, retention and terminal
bounds. Demand response is lossless shifting: a causal energy backlog records
deferred demand and must be repaid at the terminal period.

The network constrains transport capacities but does not impose Kirchhoff's
laws or contingency security. Commitment is aggregated by zone, so equipment
counts and graph nodes are not counts of independent commitment binaries.
Appendix~\ref{app:physical} gives the operational equations. All methods face
this same model and the same independent output checks.

\subsection{From a scenario to a heterogeneous temporal graph}
The scenario becomes an attributed graph with system, region, zone, asset
and weather nodes. Relations express containment, weather influence,
transmission and temporal couplings. Figures~\ref{fig:graph_spatial}
and~\ref{fig:graph_temporal} illustrate the construction. No solution or
reference objective is an input to the graph.

The hierarchical temporal encoder aggregates spatial information from assets
to zones, regions and the system, models temporal dependence, and distributes
context back to the zone sequence. Frozen training means and standard
deviations normalize the inputs; they are not refitted on evaluation data.
Online graph construction and embedding are timed. The graph supplies
descriptive context; the optimization model and validator enforce physical
constraints. A temporal edge is not itself a certificate of energy balance
or a commitment transition.
\input{paper2/figures/scenario_graphs}

\subsection{Compared methods and shared reconstruction}
Table~\ref{tab:methods} defines the actual search and reconstruction mechanisms.
HiGHS is the open-source direct baseline; its components include the parallel
dual simplex method \citep{huangfu2018}. Rolling horizon constructs a proposal
from physical inputs with a six-period look-ahead, then invokes the same repair
engine as the learned methods. It is not a solver-free method.

\begin{table}[htbp]\centering\small
\begin{tabular}{p{.17\linewidth}p{.42\linewidth}p{.29\linewidth}}\toprule
Method & Proposal or search & Online resources\\\midrule
HiGHS & Direct corrected MILP & 16 solver threads\\
RH & Rolling-horizon strategy, then repair & One repair, 16 threads\\
JEPA & Decode one predicted latent anchor, then repair & GPU; one repair, 16 threads\\
Pipeline & Sixteen energy-guided latent candidates, then repair & GPU; 16 concurrent repairs, one thread each\\
No-gradient control & Same sampler, energy gradient set to zero & Same envelope as pipeline\\
Gurobi & Independent direct solution of the same MILP & 16 solver threads\\\bottomrule
\end{tabular}
\caption{Methods and resource allocation. The no-gradient control and Gurobi belong to the prepared extensions. Every accepted output undergoes the same independent physical validation.}\label{tab:methods}
\end{table}

A strategy has eight channels per zone and time,
\begin{equation}
 u_{zt}=[b^{\rm ch},b^{\rm dis},p^{\rm ch},p^{\rm dis},d^{\rm DR},v^{\rm th},u^{\rm th},i^{\rm imp}]_{zt}.
 \label{eq:strategy}
\end{equation}
Storage channels indicate positive activity and allow idle periods. The import
label is a system binary broadcast over zones. Only the first seven channels
enter repair targets; import/export remains a reconstruction decision.
An autoencoder maps reference strategies to 32 tokens of dimension 192.
A scenario-to-latent predictor supplies an anchor without observing the
reference strategy. This follows the broad idea of predicting representations,
while differing from image-based self-supervised JEPA in its inputs, targets
and losses \citep{assran2023}.

JEPA decodes the anchor once. The pipeline initializes sixteen latents around
the anchor or from a learned prior, then applies gradient-and-noise updates
with drift toward the anchor. Energy-based modeling and Langevin updates
motivate the construction \citep{du2019,welling2011}; twenty finite updates do
not establish equilibrium sampling. Both methods use the same frozen feature
thresholds. Appendix~\ref{app:learned} specifies architectures, training
objectives, thresholds and update equations.

\input{paper2/figures/method_chain}

The repair cascade first fixes supplied targets, then releases selected
decisions under deviation penalties and flip limits. Original integer domains
and structural fixes are retained: later stages may solve MILPs. Among
independently valid candidates completed before the computation cutoff, the
pipeline returns the lowest physical cost, with candidate index breaking ties.
It does not retrospectively select another candidate to satisfy an ENS cap.
All child processes must still exit before the global deadline.

\subsection{What the matched energy ablation identifies}
The prepared control removes only the energy-gradient term. Candidate count,
models, initialization, random streams, noise, drift, clipping, update count,
thresholds and repair resources are matched. Both variants receive the same
deadline; removing gradient computation can leave more time for repair.
Their difference measures the operational contribution of that term under
this resource envelope. It does not establish a universal EBM advantage or
equate realized computation time. The one-candidate comparison alone cannot
answer this question.

\section{Experimental setup}
\subsection{Completed cohort and computational protocol}
The completed campaign contains 600 constructed scenarios, 100 per family:
low, medium and high criticality; renewable generation with batteries;
network-focused systems; and combinatorially constrained systems. The families
combine different generator and selection regimes, rather than forming six
ordered severity levels. Historical scenarios have 24 hourly periods.
There are 16,800 primary executions over seven independent budgets and 1,680
additional executions from two extra repetitions on 30 preselected scenarios.
The main analysis uses five budgets from 5 to 120 seconds; 1/2-second records
remain part of the historical record.

Execution uses one machine and an envelope of 16 CPU solver threads; learned
inference uses an RTX 4050. Static preload includes libraries, model weights
and repair workers, but no evaluation scenario or embedding. The timed
interval includes reading, graph construction, inference, repair, validation,
writing and full process exit. A fixed budget-dependent reserve leaves time
for finalization. Forced termination, late output and a still-running child
are failures, even if a valid incumbent existed earlier. The reserve remains
unchanged across the prepared extensions.

The validator reconstructs balances and costs from dispatches independently
of solver-reported objectives. Absolute physical and integrality tolerances
are $10^{-3}$ and $10^{-5}$. All 600 common references are valid; 383 are
certified optimal, while the remainder are validated incumbents from a
600-second HiGHS solve. Reference computation and training are excluded from
online timing and reported separately. A small reference-relative gap or zero
ENS is not a security or robustness guarantee.

\subsection{Outcomes and statistical comparisons}
Let $F_{smb}$ indicate valid output and complete timely exit for scenario $s$,
method $m$ and allocation $b$. Delivered unserved energy is
$E_{smb}=\sum_{z,t}P^{\rm uns}_{smbzt}\Delta t$. With reference $C_s^{\rm ref}$,
\begin{equation}
 Q_{smb}(\tau)=\begin{cases}
 \ind\{(C_{smb}-C_s^{\rm ref})/|C_s^{\rm ref}|\leq\tau\},&|C_s^{\rm ref}|>1\ {
m EUR},\\
 \ind\{C_{smb}-C_s^{\rm ref}\leq1\ {
m EUR}\},&\text{otherwise}.
 \end{cases}
\end{equation}
The primary joint endpoint is
\begin{equation}
 Y_{smb}=F_{smb}Q_{smb}(0.05)\ind\{E_{smb}\leq10^{-3}\ {\rm MWh}\}.
 \label{eq:joint}
\end{equation}
An observed failure has outcome zero without assigning fictitious cost or ENS.
Missing trials or necessary measurements remain unknown. Disabled conditions
do not require the corresponding measurement. Secondary endpoints separate
validity, continuity and economic quality; sensitivity uses 1/10\% cost
thresholds and ENS caps of 1/10/100 MWh. Absolute cost differences and ENS
relative to demand complement percentage comparisons.

Coverage retains observed failures. Conditional cost and time comparisons
use common acceptable supports and report their sizes. Pointwise 95\%
intervals use 2,000 paired bootstrap samples, resampling scenario bases within
families \citep{efron1979}. Variants, deadlines and repetitions of a base remain
together. Repetitions do not increase the number of independent scenarios.
The completed cohort was inspected during development, so its contrasts are
exploratory.

On the independent test, 25 primary contrasts are fixed in advance: RH, JEPA,
pipeline and Gurobi against HiGHS, plus pipeline against the no-gradient
control, at each deadline. Exact paired binary tests use discordant outcomes,
with Holm correction over all 25 comparisons. Missing comparisons cannot
become significant and do not reduce the correction family. Structural and
seed-sensitivity intervals remain descriptive. These distinctions address
selection-induced optimism \citep{cawley2010}.

\subsection{Independent and transfer extensions}
\label{sec:extensions}
\begin{table}[htbp]\centering\small
\begin{tabular}{p{.24\linewidth}p{.40\linewidth}p{.24\linewidth}}\toprule
Block & Sampling and separation & Evidence status\\\midrule
Historical comparison & 600 scenarios; six families; 24 h & Complete\\
Independent comparison & 120 development and 120 final cases; 20 per family in each & Prepared\\
Structural controls & 30 development and 30 final bases; seven conditions per base & Prepared\\
Multi-horizon labels & 180 separate bases at 12/24/48 h; grouped train/validation assignment & Prepared\\
External direct solver & 120 fixed historical and 120 independent final cases & Prepared\\\bottomrule
\end{tabular}
\caption{Distinct evidence blocks. New cases share family parameter spaces but not the historical selection procedure, so they are not an unqualified IID replication. Final outcomes remain unavailable until the design and models are locked.}\label{tab:cohorts}
\end{table}

Structural conditions comprise a 24-hour control, duplicated system within
300 zones, recabled network, modified thermal/storage capacities, consistent
zone permutation, and 48/72-hour extensions. Permutation preserves the physical
problem up to relabeling, including the trade anchor. Duplication changes both
size and structure. Horizon extensions preserve the first day and apply
declared daily demand/renewable scalings; they test controlled synthetic shift,
not independent forecasting errors. Transformations act on realized physical
data shared by all methods.

The historical decoder is limited to 24 periods. Longer-horizon cells are
inapplicable, not failures. A separate autoencoder, predictor and energy model
use content-conditioned queries and sinusoidal time positions, while reusing
the frozen graph encoder and normalization. Training uses 12/24/48-hour
instances; 72 hours is absent from training. Three fixed training seeds
quantify sensitivity without selecting a winner on the test. The primary seed
is evaluated across the control suite; the other two use 24/48/72 hours and
30/120-second budgets.

Gurobi solves the same physical MILP directly \citep{gurobi2026}; it does not
replace the repair backend. HiGHS references are preserved, so an improved
incumbent can have a negative gap. Two extra repetitions on 30 independent
final cases at 30/120 seconds assess variability of the principal methods
and matched ablation.

\section{Results}
\subsection{The service requirement changes the useful method}
\input{paper2/results/historical_table}
HiGHS has the strongest economic and joint coverage at 120 seconds. RH instead
provides more outputs without ENS, at a substantial cost premium. The pipeline
returns fewer valid outputs than JEPA but more outputs meeting the joint
criterion. These are properties of the evaluated implementations and resource
envelope, rather than general rankings of exact, heuristic and learned methods.

\begin{figure}[htbp]\centering
\includegraphics[width=\linewidth]{paper2/results/coverage.pdf}
\caption{Historical continuity, economic and joint coverage. Shading shows pointwise scenario-bootstrap intervals. Each deadline is an independent execution; these are not incumbent trajectories or time-to-target curves.}\label{fig:coverage}
\end{figure}

At 30 seconds, RH preserves service on all 100 high-criticality cases while
HiGHS does so on eight. At 120 seconds, HiGHS reaches the joint requirement
on 61 of those cases; the learned methods reach it on only one JEPA and two
pipeline cases. Conversely, on the renewable/battery family at ten seconds,
HiGHS meets the joint requirement on 46 cases and JEPA on 38. The union of
all four methods covers 59, including six cases served only by JEPA. This is
localized complementarity, not general superiority on that family.

\subsection{Faster acceptable outputs exist on restricted supports}
Figure~\ref{fig:paired} compares outputs meeting the same joint requirement
at the same 120-second allocation. JEPA can finish earlier on selected cases,
including 43 renewable/battery and 25 medium-criticality cases. Such comparisons
complement coverage; they cannot remove scenarios where JEPA fails. Nor does
fast completion under a 120-second allocation establish the same outcome
under an independently allocated shorter deadline.

\begin{figure}[htbp]\centering
\includegraphics[width=\linewidth]{paper2/results/paired_tradeoffs.pdf}
\caption{Paired differences against HiGHS at 120 seconds where both methods meet the joint criterion. Negative horizontal values denote earlier completion by the comparator; negative vertical values denote lower cost. Conditional support is stated in each panel.}\label{fig:paired}
\end{figure}

\subsection{Execution overhead and limits of component attribution}
\begin{figure}[htbp]\centering
\includegraphics[width=.94\linewidth]{paper2/results/execution_time.pdf}
\caption{Execution components and separate static preload at 120 seconds. Medians use terminal records with available timing, including failures. Different success supports prevent interpreting these medians as paired speedups. Concurrent repair times cannot be summed into wall-clock time.}\label{fig:time}
\end{figure}
The smaller pipeline coverage and its overhead caution against equating broader
candidate generation with greater usefulness. The pipeline--JEPA difference
mixes candidate count, initialization, energy guidance and repair opportunities.
Prepared instrumented runs additionally separate encoding, generation, repair
waiting, validation and child closure, and count changed decisions only where
the corresponding technology is present. This diagnostic measures repair
extent, not a causal mechanism.

\subsection{Sensitivity, repetition and independent evidence}
At 120 seconds, a retrospective perfect choice among the four outputs meets
the joint criterion on 477 cases, compared with 475 for HiGHS. The two additional
cases have about 0.121 and 0.682 MWh of ENS in the HiGHS outputs. With a 1 MWh
cap, HiGHS itself covers 477 cases. A small strict-continuity gain thus need
not survive a different operational tolerance.

On the 30 historically repeated cases at 120 seconds, the joint outcome changes
across repetitions for six HiGHS cases, three pipeline cases, one JEPA case
and no RH case. These counts do not establish deployment reliability, and
stable failure would not constitute useful reliability.
\input{paper2/results/extension_results}
\IfFileExists{paper2/results/energy_ablation.pdf}{
\begin{figure}[htbp]\centering\includegraphics[width=.94\linewidth]{paper2/results/energy_ablation.pdf}
\caption{Matched energy ablation, with development and independent test separate. Both methods use sixteen candidates and the same resource envelope.}\label{fig:ablation}\end{figure}}{}
\IfFileExists{paper2/results/generalization.pdf}{
\begin{figure}[htbp]\centering\includegraphics[width=\linewidth]{paper2/results/generalization.pdf}
\caption{Independent structural and horizon evaluation. N/A denotes an unsupported configuration, not a failure. Multi-horizon results use the predetermined primary seed.}\label{fig:generalization}\end{figure}}{}

\section{Discussion and limitations}
The completed evidence favors request-specific interpretation. RH offers rapid
continuity on several difficult families; HiGHS supplies stronger economic
coverage, especially with longer allocations; JEPA provides selected acceptable
outputs earlier and some unique short-budget successes. The pipeline has not
demonstrated a broad advantage, and the current comparison does not attribute
any improvement specifically to energy guidance. Negative ablation or transfer
outcomes would refine this conclusion rather than invalidate the question.

The study concerns deterministic, constructed systems with aggregated
commitment and a transport network. It does not evaluate AC feasibility,
contingency security, uncertain forecasts or deployment. Family selection
and size are partly confounded. Bootstrap intervals condition on the sampled
family mixture and do not correct this design. Complete pretraining membership
of the graph encoder requires its own provenance audit: disjoint supervised
labels alone do not certify that every learned stage is free of evaluation
exposure.

Reference-relative quality compares to a common benchmark, not an unknown
optimum. Certified-reference subsets and lower bounds provide complementary
information but can select easier cases. ENS is penalized in the objective
and reported separately, so cost and continuity are related but not
interchangeable. An absolute ENS tolerance has different relative importance
at different physical scales.

Training, labels and static preload are real costs. Excluding them from the
online interval defines a warm-service setting. Sixteen one-thread repairs
and one sixteen-thread solve have the same thread envelope, not identical
parallel efficiency. One machine and a common repair backend limit timing
portability. The variable-horizon extension changes architecture and training
data; its results will characterize that new chain, not isolate an
architectural improvement or demonstrate zero-shot transfer of the old stack.

\section{Conclusion}
Within the completed cohort, no method dominates all requested services.
At 120 seconds, HiGHS supplies the broadest joint economic/service coverage,
while RH supplies the broadest continuity coverage. Learned proposals offer
narrower time and coverage opportunities, and the sixteen-candidate pipeline's
overhead limits its usefulness. These observations answer the comparative
question within the tested envelope. The incremental value of energy guidance
and robustness under structural or temporal changes remain to be established
by the specified independent experiments.

\subsection*{Subsequent work}
Observed complementarity defines an opportunity for selection, not a selection
policy. Subsequent work will test whether information available before solving
can identify useful methods for a stated deadline and service requirement.
It must compare against fixed methods and simple baselines, use independent
evaluation, and charge feature computation and selection to the same deadline.
The criticality descriptors studied previously are candidate inputs; predictive
value for one solver does not establish a cross-method routing gain.

\appendix
\input{paper2/technical_appendix}
\begingroup\interlinepenalty=10000
\bibliographystyle{unsrtnat}\bibliography{paper2/references}
\endgroup
\end{document}
